### Theorem
Let $p$ and $q$ be distinct primes with $q > 2$. The number of conjugacy classes of subgroups of order $q$ in $G = \text{GL}_2(p)$ is:
$$s_q(\text{GL}_2(p)) = \frac{q+3}{2}\Delta^q_{p-1} + \Delta^q_{p+1}$$

### Proof

Let $G = \text{GL}_2(p)$ and let $V = \mathbb{F}_p^2$ be the 2-dimensional vector space over the finite field $\mathbb{F}_p$. Any subgroup $H \leqslant G$ of order $q$ is cyclic because $q$ is prime. Thus, $H = \langle M \rangle$ for some matrix $M \in G$ of order $q$. Since $q > 2$ is an odd prime and $p$ is a distinct prime, $q$ cannot divide both $p-1$ and $p+1$ (otherwise $q$ would divide their difference, $2$). This yields two mutually exclusive cases.

#### Case 1: $q \mid (p+1)$ and $q \nmid (p-1)$

We first establish that any matrix $M \in G$ of order $q$ has no eigenvalues in the base field $\mathbb{F}_p$. Suppose for contradiction that $\lambda \in \mathbb{F}_p$ is an eigenvalue of $M$, so $Mv = \lambda v$ for some non-zero $v \in V$. Since $M^q = I$, we must have $v = M^q v = \lambda^q v$, which implies $\lambda^q = 1$ in $\mathbb{F}_p$. By Lagrange's Theorem, the multiplicative order of $\lambda$ in $\mathbb{F}_p^*$ must divide $|\mathbb{F}_p^*| = p-1$. Because $q \nmid (p-1)$, the order of $\lambda$ cannot be $q$. Thus, the order of $\lambda$ is $1$, meaning $\lambda = 1$. 

If $\lambda = 1$ is an eigenvalue, the characteristic polynomial of $M$ over $\mathbb{F}_p$ must be $f(x) = (x-1)(x-d)$ for some $d \in \mathbb{F}_p$. Since $d$ is an eigenvalue of $M$ in $\mathbb{F}_p$, the same argument forces $d^q = 1$ in $\mathbb{F}_p$. Since $q \nmid (p-1)$, we must have $d = 1$, so $f(x) = (x-1)^2$. Because $M^q - I = 0$, the minimal polynomial of $M$ must divide $x^q - 1$. Since $q$ and $p$ are distinct primes, $x^q - 1$ has no multiple roots in any extension of $\mathbb{F}_p$. Thus, the minimal polynomial of $M$ has simple roots, implying $M$ is diagonalizable over its splitting field. Since its only eigenvalue is $1$ and it is diagonalizable, we must have $M = I$. This contradicts the fact that $M$ has order $q > 2$. Thus, $M$ has no eigenvalues in $\mathbb{F}_p$.

Since $M$ has no eigenvalues in $\mathbb{F}_p$, its degree-2 characteristic polynomial $f(x) = x^2 - tx + d \in \mathbb{F}_p[x]$ has no roots in $\mathbb{F}_p$ and is therefore irreducible. We can construct a quadratic field extension $K = \mathbb{F}_p[x]/\langle f(x) \rangle \cong \mathbb{F}_{p^2}$. In this field, the element $\alpha = x + \langle f(x) \rangle$ is a root of $f(x)$. Since $M^q = I$, we have $\alpha^q = 1$. Because $q \nmid (p-1)$, $\alpha \notin \mathbb{F}_p$, meaning $\alpha$ is a primitive $q$-th root of unity. 

The multiplicative group $K^*$ is cyclic of order $p^2-1$. Since $q \mid (p+1)$, we have $q \mid (p-1)(p+1) = p^2-1$. Thus, $K^*$ contains a unique cyclic subgroup of order $q$. Let $\omega \in K^*$ generate this unique subgroup. Every primitive $q$-th root of unity in $K$ must be of the form $\omega^j$ for some $1 \leqslant j < q$. 

Now, let $H_1 = \langle M_1 \rangle$ and $H_2 = \langle M_2 \rangle$ be two subgroups of order $q$ in $G$. The roots of the characteristic polynomial of $M_1$ are $\{\omega^a, \omega^{ap}\}$ for some $1 \leqslant a < q$, and the roots of the characteristic polynomial of $M_2$ are $\{\omega^b, \omega^{bp}\}$ for some $1 \leqslant b < q$. Since $q$ is prime, we can find an integer $k \in \{1, \dots, q-1\}$ such that $kb \equiv a \pmod q$. The matrix $M_2^k$ is a generator of $H_2$ and its eigenvalues are $\{(\omega^b)^k, (\omega^{bp})^k\} = \{\omega^a, \omega^{ap}\}$. Thus, $M_1$ and $M_2^k$ share the exact same irreducible characteristic polynomial $h(x) = x^2 - Tx + D$.

For any matrix $M \in G$ with irreducible characteristic polynomial $x^2 - Tx + D$, we can choose any non-zero vector $v \in V$. Since $M$ has no eigenvalues, $Mv$ is not a multiple of $v$, so $\{v, Mv\}$ is a basis of $V$. Under this basis, because $M(Mv) = M^2 v = TMv - Dv$, the matrix representation of the linear transformation is the companion matrix:
$$C = \begin{pmatrix} 0 & -D \\ 1 & T \end{pmatrix}$$
Consequently, both $M_1$ and $M_2^k$ are conjugate to the same companion matrix $C$. By transitivity, they are conjugate to each other, so there exists some $P \in G$ such that $P^{-1} M_1 P = M_2^k$. This implies $P^{-1} H_1 P = H_2$, proving that any two subgroups of order $q$ are conjugate. Thus, there is exactly $1$ conjugacy class of subgroups of order $q$ when $q \mid (p+1)$.

#### Case 2: $q \mid (p-1)$ and $q \nmid (p+1)$

We first show that any subgroup $H = \langle M \rangle$ of order $q$ must have eigenvalues in $\mathbb{F}_p$. If $M$ had no eigenvalues in $\mathbb{F}_p$, then by the proof in Case 1, its eigenvalues would lie in $\mathbb{F}_{p^2} \setminus \mathbb{F}_p$ and have order $q$. This would imply $q \mid (p^2-1)$ and $q \nmid (p-1)$, which requires $q \mid (p+1)$, contradicting our assumption that $q \nmid (p+1)$. Therefore, $M$ has eigenvalues in $\mathbb{F}_p$. 

Since the eigenvalues of $M$ lie in $\mathbb{F}_p$, $M$ is diagonalizable over $\mathbb{F}_p$. Thus, up to conjugacy, every subgroup $H$ of order $q$ is a subgroup of the diagonal subgroup $D \cong C_{p-1} \times C_{p-1}$. Let $a \in \mathbb{F}_p^*$ be an element of order $q$. Any subgroup of order $q$ in $D$ is cyclic and generated by a diagonal matrix of the form $\text{diag}(a^i, a^j)$ where $(i, j) \not\equiv (0, 0) \pmod q$. We can classify these subgroups by a parameter $\ell \in \mathbb{F}_q \cup \{\infty\}$:
* If $i \equiv 0 \pmod q$, we must have $j \not\equiv 0 \pmod q$. Rescaling the generator yields the subgroup $U_{\infty} = \langle \text{diag}(1, a) \rangle$.
* If $i \not\equiv 0 \pmod q$, we can rescale the generator by raising it to the power $i^{-1} \pmod q$, which yields a unique generator of the form $\text{diag}(a, a^{\ell})$ for some $\ell \in \mathbb{F}_q$. This gives the subgroups $U_{\ell} = \langle \text{diag}(a, a^{\ell}) \rangle$.

The normalizer of $D$ in $G$ is the monomial group $I = C_2 \ltimes D$, where the non-trivial element of $C_2$ acts by swapping the diagonal entries. Under this swap, the generator of $U_{\ell}$ becomes $\text{diag}(a^{\ell}, a)$.
* If $\ell \equiv 0$, the swap maps $U_0 = \langle \text{diag}(a, 1) \rangle$ to $\langle \text{diag}(1, a) \rangle = U_{\infty}$. Thus, $U_0$ and $U_{\infty}$ are conjugate.
* If $\ell \not\equiv 0$, we can normalize the first entry of $\text{diag}(a^{\ell}, a)$ back to $a$ by raising it to the power $\ell^{-1} \pmod q$. This yields the generator $\text{diag}(a, a^{\ell^{-1}})$. Thus, the swap maps $U_{\ell}$ to $U_{\ell^{-1}}$.

Therefore, two diagonal subgroups $U_{\ell}$ and $U_k$ are conjugate in $G$ if and only if $\ell$ and $k$ belong to the same orbit of the action $\ell \mapsto \ell^{-1}$ on the parameter space $\mathbb{F}_q \cup \{\infty\}$. The orbits under this action are:
1. $\{0, \infty\}$, which forms $1$ orbit of size 2.
2. $\{1\}$, which forms $1$ orbit of size 1 (since $1^{-1} \equiv 1$).
3. $\{-1\}$ (or $q-1$), which forms $1$ orbit of size 1 (since $(-1)^{-1} \equiv -1$).
4. The remaining $q-3$ elements of $\mathbb{F}_q^* \setminus \{1, -1\}$ do not equal their own inverses. Because $q > 2$ is prime, these elements pair up into exactly $\frac{q-3}{2}$ orbits of size 2.

The total number of orbits is:
$$1 + 1 + 1 + \frac{q-3}{2} = 3 + \frac{q-3}{2} = \frac{q+3}{2}$$

Each orbit corresponds to exactly $1$ conjugacy class of subgroups of order $q$. Thus, there are exactly $\frac{q+3}{2}$ conjugacy classes of subgroups of order $q$ when $q \mid (p-1)$. 

By combining the two cases, we obtain the desired formula:
$$s_q(\text{GL}_2(p)) = \frac{q+3}{2}\Delta^q_{p-1} + \Delta^q_{p+1}$$

$\blacksquare$